\documentclass[11pt]{article}
\usepackage[a4paper, top=18mm, bottom=18mm, left=20mm, right=20mm]{geometry}
\usepackage[T2A]{fontenc}
\usepackage[utf8]{inputenc}
\usepackage[ukrainian]{babel}
\usepackage{graphicx}
\usepackage[export]{adjustbox}
\usepackage{amsmath}
\usepackage{amssymb}
\setlength{\parindent}{0pt}
\setlength{\parskip}{4pt}
\pagestyle{empty}
\begin{document}
%begin
{\large\textbf{Контрольна робота}}

\textbf{Варіант \No\,1}\hfill \textit{ПІБ:}\,\underline{\hspace{60mm}}
\vspace{4mm}

%beginitem
1. %goodpath:Scripts/2018/Mod2018_1/03-good.txt
Відмітити все, що є коректним оператором С++:
( 1 ) double n.m;

( 2 ) if a<b a=0;

( 3 ) {int x,y;}

( 4 ) double(2)/87

%enditem
%beginitem
2. 
try:
    res = 9**-1.0*False
    if isinstance(res, bool):
        print(f'bool;{res}')
    elif isinstance(res, int):
        print(f'int;{res}')
    elif isinstance(res, float):
        print(f'float;{res:.2f}')
    else:
        print('???')
except:
    print('Z')

%enditem
%beginitem
3. Що буде виведено за виконання фрагменту коду?
%code
if (5 != 5)
    cout << "f";
else
    cout << "s";
%code


%enditem
%beginitem
4. %goodpath:Scripts/2019/Mod2019_1/Lex/01-good.txt
Відмітити все, що є однією правильною лексемою:
( 1 ) 5_w

( 2 ) 0x17

( 3 ) 'c'

( 4 ) "1,3,2"

%enditem
%beginitem
5. 
code
d = 0
k = 1

class D:
    k = 2
    
    def __init__(self):
global k        self.k = 3
        k = 4
        D.k = 5

obj = D()
print(d, k, D.k, obj.k, sep=' ')
code

%enditem
%beginitem
6. 
try:
    print(5, end=';')
    print(2j!=5j, end=';')
    print(0, end=';')
except TypeError:
    print(9, end=';')
except ZeroDivisionError:
    print(6, end=';')
except:
    print('ERR', end=';')
else:
    print(8, end=';')
finally:
    print(1)

%enditem
%beginitem
7. Що буде виведено за виконання фрагменту коду?

%code
def h(a,b):
    c=84
    if b!=-3:
        c=5
    elif a>-2:
         return 2
    else: 
        return 0
    return c

print(h(3,3))
%code

%enditem
%beginitem
8. %goodpath:Scripts/2019/Mod2019_1/Lex/03-good.txt
Відмітити все, що є коректним оператором С++:
( 1 ) x not is y

( 2 ) print(end="1","qwe")

( 3 ) 3**-1

( 4 ) x-=2

%enditem
%beginitem
9. 
Записати ВИРАЗ, значенням якого є словник, ключами якого є цілі числа від 12 до 2012 (включно), що не діляться на жодне з чисел 2, 3, а ключам відповідають їх квадратні корені 
%enditem
%beginitem
10. Що буде виведено за виконання фрагменту коду?

%code
a,b,c=5,2,0
def h(b):
global c    a*=3
    b=4
    c=2
    return a+b+c

a,b,c=9,6,8
print(h(b),a,b,c)
%code

%enditem










%end
\newpage

\newpage

%begin
{\large\textbf{Контрольна робота}}

\textbf{Варіант \No\,2}\hfill \textit{ПІБ:}\,\underline{\hspace{60mm}}
\vspace{4mm}

%beginitem
1. 

print('R', end=' ')
f=['0','1','2','3','4']
print(f.pop(0), end=' ')
print(f)


%enditem
%beginitem
2. 
Написати програму, що запитує у користувача значення дійсних параметрів $x$ та $y$, обчислює для них значення виразу та повiдомляє введенi данi та результат обчислення.
$$\frac{\log_{2} x - 44}{(\tg x - 0.3) (y - 99)}$$ 


%enditem
%beginitem
3. %goodpath:Scripts/2018/Mod2018_1/03-good.txt
Відмітити все, що є коректним оператором С++:
( 1 ) double n.m;

( 2 ) if a<b a=0;

( 3 ) {int x,y;}

( 4 ) double(2)/87

%enditem
%beginitem
4. Що буде виведено за виконання фрагменту коду?

%code
#include <iostream>

int h(int a, int b){
    int c = 84;
    if (b != -3)
        c = 5;
    else if (a > -2)
         return 2;
    else 
        return 0;
    return c;
}

int main(){
    std::cout << h(3, 3);
    return 0;
}
%code

%enditem
%beginitem
5. 
try:
    res = "False"!=True
    if isinstance(res, bool):
        print(f'bool;{res}')
    elif isinstance(res, int):
        print(f'int;{res}')
    elif isinstance(res, float):
        print(f'float;{res:.2f}')
    else:
        print('???')
except:
    print('Z')

%enditem
%beginitem
6. Що буде виведено за виконання фрагменту коду?

%code
f = ('a','b','c','d','e',0,1,2,3)
print(f[-1])
%code


%enditem
%beginitem
7. %goodpath:Scripts/2019/Mod2019_1/Lex/03-good.txt
Відмітити все, що є коректним оператором С++:
( 1 ) x not is y

( 2 ) print(end="1","qwe")

( 3 ) 3**-1

( 4 ) x-=2

%enditem
%beginitem
8. %goodpath:Scripts/2018/Mod2018_1/01-good.txt
Відмітити все, що є однією правильною лексемою:
( 1 ) 5_w

( 2 ) (1)

( 3 ) )

( 4 ) 0.5

%enditem
%beginitem
9. 
a = 5
b = 2
c = 0

def h(b):
global c    a *= 3
    b = 4
    c = 2
    return a + b + c

a = 9
b = 6
c = 8

try:
    print(h(b), a, b, c, sep=';')
except:
    print('Z', a, b, c, sep=';')

%enditem
%beginitem
10. Що буде виведено за виконання фрагменту коду?

%code
print('R', end=' ')
f=[0,1,2,3,4,5,6,7,8,9]
print(f[:-3:-3])
%code


%enditem










%end
\newpage

\newpage

%begin
{\large\textbf{Контрольна робота}}

\textbf{Варіант \No\,3}\hfill \textit{ПІБ:}\,\underline{\hspace{60mm}}
\vspace{4mm}

%beginitem
1. 

def h(a,b):
    c=84
    if b!=-3:
        c=5
    elif a>-2:
         return 2
    else: 
        return 0
    return c

print(h(3,3))
%enditem
%beginitem
2. %goodpath:Scripts/2019/Mod2019_1/Lex/01-good.txt
Відмітити все, що є однією правильною лексемою:
( 1 ) 5_w

( 2 ) 0x17

( 3 ) 'c'

( 4 ) "1,3,2"

%enditem
%beginitem
3. 
try:
    t = 5
    while t<=7:
        t += 2
    print(t, end=';')
except:
    print('ERR')

%enditem
%beginitem
4. 
Записати ВИРАЗ, значенням якого є множина, що складається з кубівцілих чисел від 12 до 2012 (включно), що не діляться на жодне з чисел 2, 3.
%enditem
%beginitem
5. 
Змінні \kw{next} та \kw{prev} вузла списку позначають наступний та попередній за ним вузол відповідно. Починаючи з голови, у список записано послідовні цілі числа від~$-21$ до~$20$. Нехай змінна \kw{p} позначає вузол, що зберігає значення~$3$.
Яке число зберігається у вузлі \kw{p.next.next.next.next.next} ?

%enditem
%beginitem
6. %goodpath:Scripts/2019/Mod2019_1/Lex/03-good.txt
Відмітити все, що є коректним оператором С++:
( 1 ) x not is y

( 2 ) print(end="1","qwe")

( 3 ) 3**-1

( 4 ) x-=2

%enditem
%beginitem
7. 
def f(n): print("f", end=""); return n!=-2
   print(f(3) or f(-6))
%enditem
%beginitem
8. 
for f in range(-3,-3,-3):
    if f>=8:
        break
        print(f,end=' ')
    if f>=8:
        break
else:
    print(f, end=' ')
print(f, end=' ')

%enditem
%beginitem
9. %goodpath:Scripts/2018/Mod2018_1/01-good.txt
Відмітити все, що є однією правильною лексемою:
( 1 ) 5_w

( 2 ) (1)

( 3 ) )

( 4 ) 0.5

%enditem
%beginitem
10. Що буде виведено за виконання фрагменту коду?

%code
for f in range(2,2+4,3):
    if f>=8:
        pass
    print(f, end=' ')
    f=2
    if f>=8:
        break
else:
    print(f, end=' ')
print(f, end=' ')
%code

%enditem










%end
\newpage

\newpage

%begin
{\large\textbf{Контрольна робота}}

\textbf{Варіант \No\,4}\hfill \textit{ПІБ:}\,\underline{\hspace{60mm}}
\vspace{4mm}

%beginitem
1. Що буде виведено за виконання фрагменту коду?

%code
cout << (false != 7 > 4 < 2);
%code

%enditem
%beginitem
2. 
Записати ВИРАЗ, значенням якого є множина, що складається з кубівцілих чисел від 12 до 2012 (включно), що не діляться на жодне з чисел 2, 3.
%enditem
%beginitem
3. %goodpath:Scripts/2019/Mod2019_1/Lex/03-good.txt
Відмітити все, що є коректним оператором С++:
( 1 ) x not is y

( 2 ) print(end="1","qwe")

( 3 ) 3**-1

( 4 ) x-=2

%enditem
%beginitem
4. Що буде виведено за виконання фрагменту коду?

%code
print('R', end=' ')
f=[0,1,2,3,4,5,6,7,8,9]
print(f[:-3:-3])
%code


%enditem
%beginitem
5. %goodpath:Scripts/2018/Mod2018_1/01-good.txt
Відмітити все, що є однією правильною лексемою:
( 1 ) 5_w

( 2 ) (1)

( 3 ) )

( 4 ) 0.5

%enditem
%beginitem
6. Що буде виведено за виконання фрагменту коду?

%code
a,b,c=5,2,0
def h(b):
    a=3
    b=4
    c=2
    return a+b+c

a,b,c=9,6,8
print(h(b),a,b,c)
%code

%enditem
%beginitem
7. %goodpath:Scripts/2019/Mod2019_1/Lex/01-good.txt
Відмітити все, що є однією правильною лексемою:
( 1 ) 5_w

( 2 ) 0x17

( 3 ) 'c'

( 4 ) "1,3,2"

%enditem
%beginitem
8. 
Значенням змінної \kw{a} є цілочисловий масив типу $numpy.ndarray$ розмірів $5{\times}5$. На скільки збільшиться сума елементів масиву \kw{a} після виконання
\begin{verbatim}
a.T[2] += 3
\end{verbatim}


Якщо код некоректний, то в якості відповіді зазначити ERROR.



%enditem
%beginitem
9. 
try:
    print(5, end=';')
    print(2j!=5j, end=';')
    print(0, end=';')
except TypeError:
    print(9, end=';')
except ZeroDivisionError:
    print(6, end=';')
except:
    print('Z', end=';')
else:
    print(8, end=';')
finally:
    print(1)

%enditem
%beginitem
10. Що буде виведено за виконання фрагменту коду?

%code
#include <iostream>
using namespace std;

int h(int a, int b){
    int c = 84;
    if (b != -3)
        c = 5;
    else if (a > -2)
         return 2;
    else 
        return 0;
    return c;
}

int main(){
    cout << h(3, 3);
    return 0;
}
%code

%enditem










%end
\newpage
\end{document}


